\chapter{外部意见的交叉核验}\label{app:crosscheck}

本书的框架在写作过程中接受了两条独立的外部意见通道（Grok 与 Gemini，2026-10-31），
以及一份真实数据的检验。这一章记录：他们说了什么、哪些被采纳、哪些被数据推翻、
哪些被本书修正。保留这段记录，是因为\textbf{"我参考了谁"和"我为什么采信"本身
就是研究可信度的一部分}。

\section{两条通道的分工}

\begin{center}
\begin{tabular}{@{}p{0.16\textwidth}p{0.34\textwidth}p{0.40\textwidth}@{}}
\toprule
通道 & 强项 & 本书用它问什么 \\
\midrule
Grok & 实务细节、数据源现状、参数区间 & 数据源可用性、因子清单、成本参数、模型预期 \\
Gemini & 数学推导、文献细节、常见误用 & DSR/CSCV 公式、purge/embargo、NW 滞后阶数、FDR \\
\bottomrule
\end{tabular}
\end{center}

\section{被采纳的意见}

\subsection{关于因子方向}

两条通道都指出：\textbf{A 股日频上中长期动量偏弱甚至为负，短期反转主导}。
Grok 的原话是"经典中长期动量在 A 股长期偏弱甚至为负"。

本书的实测完全支持这一点：20 日动量的 Rank IC 约为 $-0.05$ 到 $-0.08$。
这条意见直接改变了第~\ref{ch:factor} 章的写法——书中明确写出"动量方向必须实证，
不能照搬教科书"。

\subsection{关于印花税}

Grok 指出 2023-08-28 起证券交易印花税由 $0.1\%$ 降至 $0.05\%$（单边卖出）。
本书核对后确认，并据此把双边显性成本估算从约 $0.13\%$ 修正为约 $0.08\%$--$0.11\%$。

\subsection{关于 IC 的年化口径}

Gemini 特别强调：日频算 IC 但预测 $h$ 日收益时，年化乘数应为 $\sqrt{252/h}$
而不是 $\sqrt{252}$，否则会虚增 $2$--$4$ 倍。本书第~\ref{ch:eval} 章把这个陷阱
单独列为一节，并在实现中采用 $\sqrt{252/5}$。

\subsection{关于 Purge 与 Embargo 的区间}

Gemini 给出了完整形式：训练样本需剔除
$[t_{\text{test,start}} - h - w,\ t_{\text{test,stop}} + h + w]$，
其中 $h$ 是标签持有期、$w$ 是特征回看窗。本书的 \texttt{aq/cv.py} 采用其工程简化版
（左侧 purge 到 $t_0-h$，右侧 embargo 到 $t_1+w$），并在第~\ref{ch:overfit} 章说明了理由。

\subsection{关于多重检验}

Gemini 给出了 Benjamini-Hochberg 的完整步骤，并指出一个关键细节：
\textbf{必须从大到小找最大的 $k^*$，而不是从小到大遇到第一个不满足就停}。
本书第~\ref{ch:eval} 章的代码实现遵循这一点。

\section{被数据修正的意见}

\subsection{"中性化后 IC 一定更高"}

Grok 在初稿中倾向于"中性化能提升因子质量"。本书的实测显示：\textbf{并非如此}。
部分因子的中性化后 IC 明显低于原始 IC（例如规模因子对市值中性化后退化为零）。
第~\ref{ch:prep} 章因此专门写了"中性化不总是提升 IC"一节。

\subsection{"模型能带来显著提升"}

Grok 给出的预期是"GBDT 比线性提升 0.01--0.02 的 IC"。本书的实测显示，
在严格 Purged 验证下，LightGBM 相对等权基线的提升\textbf{幅度有限且不稳定}
（见第~\ref{ch:result} 章的协议表）。这与"模型是次要的、因子与验证才是主要的"
这一判断一致。

\subsection{数据源的"可用性"}

两条通道都给出了数据源清单，但没有一条预见到当天的实际情况：
东方财富在约 50 次请求后被 IP 封锁，腾讯在约 200 次后触发 WAF，
baostock 并发即被限速。最终可用的是新浪。
\textbf{这说明"数据源对比表"这类静态知识，在工程实践中的保质期极短。}
第~\ref{ch:data} 章因此把这段排雷过程完整保留下来，而不是给一张漂亮的对比表。

\section{两条通道之间的分歧}

\begin{center}
\begin{tabular}{@{}p{0.16\textwidth}p{0.30\textwidth}p{0.30\textwidth}p{0.14\textwidth}@{}}
\toprule
议题 & Grok 的观点 & Gemini 的观点 & 本书处理 \\
\midrule
三重障碍法 & 认为日频选股不必优先使用 & 未直接评价，但强调重叠标签问题 & 采纳 Grok：日频不用 \\
深度模型 & 认为 2026 年值得一试 & 未涉及 & 保留但降级为"可选" \\
PBO 的解读 & 未深入 & 强调不能判断绝对收益 & 采纳 Gemini，单列一节 \\
换手与成本 & 给出 15--40 bp 区间 & 未涉及 & 采纳 Grok，取 30 bp \\
\bottomrule
\end{tabular}
\end{center}

\section{核验原则}

本书对外部意见的处理遵循三条原则：

\begin{enumerate}
  \item \textbf{可核验的事实回原始来源}。印花税率、指数编制规则、法规条文，
        一律回官方口径，不采信二手转述。
  \item \textbf{可计算的结论自己算一遍}。DSR、PBO、ICIR 的口径，全部用本书的实现重算，
        不直接引用别人的数字。
  \item \textbf{有分歧的地方做实验}。因子方向、模型提升幅度这类问题，
        用本书的数据跑一遍，以实测为准。
\end{enumerate}

\begin{warnbox}[一个必须承认的局限]
外部意见通道（以及本书作者）都可能产生幻觉。本章记录的是"参考了什么"，
而不是"因为谁说了所以对"。所有进入正文的结论，都在本书的数据上重新验证过；
未能验证的，正文中会明确标注为"经验判断"或"待验证"。
\end{warnbox}
